\scanpage{68}
\[
(1)\qquad\Pi_1(X)\longrightarrow\widehat\Pi_1(X)
\]
which induces a bijection on the set of objects \inserted{or ``vertices''} -- identified with $X=X(\mathbb C)$ -- and which \deleted{defines} for every $x\in X=X(\mathbb C)$, gives the group homomorphism
\[
(2)\qquad
\begin{array}{ccc}
\operatorname{Aut}_{\Pi_1(X)}(x)&\longrightarrow&\operatorname{Aut}_{\widehat\Pi_1(X)}(x)\\
\Vert\ \text{definition}&&\Vert\ \text{definition}\\
\pi_1(X,x)&&\widehat\pi_1(X_{\mathrm{sch}},x)
\end{array}
\]
allowing one to identify the algebraic $\pi_1$ with the profinite completion of the ``topological'' or ``transcendental'' $\pi_1$ : \unsure{b}
\[
(3)\qquad\pi_1(X_{\mathrm{sch}},x)\simeq\widehat\pi_1(X,x).
\]
Likewise, for $x,y\in X$, the map
\[
\operatorname{Hom}_{\Pi_1(X)}(x,y)\longrightarrow\operatorname{Hom}_{\widehat\Pi_1(X)\deleted{\unsure{y}}}(x,y)
\]
\margin{We shall identify an element of $\operatorname{Hom}_{\Pi_1}(x,y)$ with its image in $\operatorname{Hom}_{\widehat\Pi_1}(x,y)$, without worrying whether the map is injective or not.}
allows one to identify the second set, as a bi-torsor over $(\widehat\pi_1(X,y),\widehat\pi_1(X,x))$ with the bitorsor obtained \deleted{Hom} from the first set by extension of the groups of operators by
\[
\pi_1(X,y)\longrightarrow\widehat\pi_1(X,y),\qquad\pi_1(X,x)\longrightarrow\widehat\pi_1(X,x).
\]
(this pair of homomorphisms satisfying \inserted{indeed} the condition required for this extension to make sense\ldots).
\margin{namely (in terms of $(G,H),(G',H')$ and $G\to G',H\to H'$ : there exists an outer isomorphism $G'\xrightarrow{\sim}H'$ making commutative the diagram \deleted{\unsure{of bitorsors}}
\diagram{\begin{tikzcd}[ampersand replacement=\&,column sep=large,row sep=large]
G\arrow[r,"\sim", "\substack{\text{defined by the}\\\text{bitorsor }(G,H)}"']\arrow[d]\&H\arrow[d]\\
G'\arrow[r,"\sim"]\&H'
\end{tikzcd}})}
Let
\[
(5)\qquad\mathbf{\Gamma}_{\deleted{K}}=\mathbf{\Gamma}_K=\operatorname{Aut}_{K\text{-cor}}(\mathbb C),
\]
\margin{67}
